\Needspace{10\baselineskip}
\subsection{Hardness consequences}
\label{subsec:bounded-degree-consequences}

Fix a target factor $r$ below a dense hardness threshold, and choose
a dense factor $r'>r$. Write $N$ for the dense graph's order and
$Y$ for its YES upper bound. It is linear
in $N$ for treedepth and treewidth and quadratic for total depth.
For $\tau>0$, the sparsification losses $\tau N$ and $\tau N^2$ can therefore be made
at most $\delta Y$ by choosing $\tau$ small. Taking
$0<\delta<(r'-r)/(r+1)$, Lemma~\ref{lem:gap-preservation} transfers
the factor-$r$ hardness.
We verify the factor conditions below, using the following observation
for the SSE construction.
In these applications we choose connected replacements using
Lemma~\ref{lem:factor-expanders}. The biclique side-size ratios,
when bicliques occur, are bounded as checked below. Although the
general sparsification theorems do not require balance, these
connected replacements preserve the connected components of the
dense graph: each replaced factor still connects its entire support.

The transfer steps are summarized below. Here $H$ is a dense graph,
$H'$ its sparse replacement, $R$ its regular extension, and
$R_\ell$ replaces each edge of $R$ by a path of fixed even length
$\ell$. Sparsification and regularization apply to all three objectives;
subdivision applies to treedepth and treewidth, and to total depth under SSE.
The guarantees are Theorems~\ref{thm:td-sparsification}--\ref{thm:total-depth-sparsification},
Proposition~\ref{prop:regularization}, and
Corollary~\ref{cor:subdivision-hardness}, respectively:
\[
 H\xrightarrow{\text{sparsify}}H'
 \xrightarrow{\text{regularize}}R
 \xrightarrow{\text{subdivide}}R_\ell.
\]

\begin{lemma}[Clique factors of a line graph]\label{lem:linegraph-factors}
Let $G$ be an unweighted host multigraph. Identify each edge
occurrence $e\in E(G)$ with its vertex in $\mathcal L(G)$, and put
$U_v=\{e\in E(G):v\text{ is an endpoint of }e\}$ for $v\in V(G)$.
The nonempty sets
$U_v$ give clique factors for $\mathcal L(G)$, and each line-graph
vertex belongs to at most two factors.
\end{lemma}

\begin{proof}
Two line-graph vertices are adjacent if and only if their host edges
share an endpoint, so
$E(\mathcal L(G))=\bigcup_{v\in V(G)}\binom{U_v}{2}$.
An edge has at most two endpoints. A loop contributes one vertex
to the factor at its endpoint.
\end{proof}

We may require arbitrarily large sources: disjoint copies preserve the satisfiable fraction
of a parity system and the independence fraction and maximum degree
of an independent-set source. A relative independent-set
approximation on the copies transfers to the original by choosing
the copy with the largest selected set. For SSE, parallel copies
of every host edge preserve normalized volumes and normalized cut sizes and
multiply the line-graph order.

In all SSE cases, let $c>0$ and $\varepsilon_0>0$ be the constants
in Lemma~\ref{lem:sse-unweighted-gap}. Apply it with sufficiently
small source error $0<\eta<\varepsilon_0$ and
$q=\lceil2/\eta\rceil$, as in Section~\ref{subsec:sse-dense-gap}.
Let $G^*$ denote its host, put $H=\mathcal L(G^*)$, and write
$N=|E(G^*)|=|V(H)|$, counting edge multiplicity, and
$\beta=(c/2)\sqrt\eta\le1/4$.
The line graph has YES height bound $4\eta N$ and NO balanced-separator
bound $\beta N$, the latter by
Lemma~\ref{lem:linegraph-separators} with $\rho=1/2$.
Lemma~\ref{lem:linegraph-factors} gives $d_1=2$, $d_0=0$, and no
biclique factors, regardless of the host degree. Sparsification
replaces these host-incidence cliques directly; the sparse output
need not itself be a line graph. Since $\beta/\eta$ grows
without bound as $\eta$ decreases, any fixed target gap can be
obtained before sparsification.

\Needspace{7\baselineskip}
\begin{corollary}[Bounded-degree treedepth hardness]
\label{cor:td-bounded-degree}
\label{thm:np-bounded-degree}
\label{thm:ugc-bounded-degree}
\label{thm:sse-bounded-degree}
For each fixed $r\ge1$, approximating treedepth within factor $r$
on simple unweighted graphs of bounded maximum degree is
NP-hard if $r<18/17$, Unique Games-hard if $r<4/3$, and
SSE-hard for every fixed $r$.
The degree bound depends only on the excluded factor. All reductions
are deterministic.
\end{corollary}

\begin{proof}
Apply Theorem~\ref{thm:td-sparsification} and
Lemma~\ref{lem:gap-preservation} in each case.

\emph{NP-hardness.}
Retain every edge of $G_F$, whose maximum degree is at most
$3b_\epsilon$ by Lemma~\ref{lem:xor-cover}. After omitting
$z_A,z_B,z_C$, replace the three anchor cliques and the three
anchor--part bicliques. Each vertex lies in at most two factor
supports, and each biclique has both sides of size $3m$.
Thus $d_1=2$, $d_0=3b_\epsilon$, and $\kappa=1$.
Deleting the three designated vertices costs at most $3$ in treedepth.
The remaining order is $18m$, so sparsification costs at most $18\tau m$.
The dense gap has positive linear slack; choose $\tau$ small and then
$m$ large. Appendix~\ref{app:transfer-td-np} gives the exact bounds
and choices.

\emph{UGC hardness.}
Fix a target factor $r<4/3$, and choose $r'$ with
$r<r'<4/3$. Take $q,\eta$ from the dense gap calculation in
Section~\ref{subsec:ugc-td-gap} for factor $r'$, and apply
Theorem~\ref{thm:bounded-degree-kr-source} with this $\eta$ in
$S_q(G)$, where $n=|V(G)|$. Replace the $q$ layer cliques by
bounded-degree expanders and retain all cross-layer source edges. These retained edges have
maximum degree at most $(q-1)d_\eta$, so $d_1=1$ and
$d_0=(q-1)d_\eta$. A layer clique gives a YES scale of at least $n$, while the graph
has $qn$ vertices. Thus the loss is at most $q\tau Y$ and can be
absorbed into the factor slack; see Appendix~\ref{app:transfer-td-ugc}.

\emph{SSE hardness.}
Use the common line-graph factors above.
Lemma~\ref{lem:balanced-by-td} converts the separator bound into
NO treedepth bound $\beta N$. Since $\beta/\eta$ is unbounded as $\eta$ tends to zero, first
choose a dense gap larger than the target factor, then make $\tau$
small enough to preserve it. The calculation is in
Appendix~\ref{app:transfer-td-sse}.

\end{proof}

The NP-hardness assertion answers the bounded-degree conjecture of
Schirrmacher, Siebertz, and
Vigny~\cite[Conjecture~1]{SchirrmacherSiebertzVigny2025}.
Weighted treedepth hardness at maximum degree ten is proved
in~\cite{DirksEtAl2025Weighted}.

\Needspace{7\baselineskip}
\begin{corollary}[Bounded-degree treewidth hardness]
\label{cor:tw-bounded-degree}
\label{thm:tw-bounded-degree}
\label{thm:tw-ugc-bounded-degree}
\label{thm:tw-sse-bounded-degree}
For each fixed $r\ge1$, approximating treewidth within factor $r$
on simple unweighted graphs of bounded maximum degree is
NP-hard if $r<12/11$, Unique Games-hard if $r<4/3$, and
SSE-hard for every fixed $r$.
The degree bound depends only on the excluded factor. All reductions
are deterministic.
\end{corollary}

\begin{proof}
Apply Theorem~\ref{thm:tw-sparsification} and
Lemma~\ref{lem:gap-preservation}.

\emph{NP-hardness.}
In $K_R(G_F)$, replace the three global cliques and retain every
edge of each source copy. Lemma~\ref{lem:xor-cover} bounds the
retained degree by $d_0=3b_\epsilon$. Each vertex lies in one
clique factor, so $d_1=1$. There are no biclique factors.
The clique replacements use $\kappa=1$ in
Lemma~\ref{lem:factor-expanders}.
The dense gap has positive slack proportional to $m$, and the
sparsification loss is $9\tau Rm$. A sufficiently small fixed $\tau$
preserves the gap; Appendix~\ref{app:transfer-tw-np} records its
coefficient and the exact choice.

\emph{UGC hardness.}
Fix a target factor $r<4/3$, and choose $r'$ with
$r<r'<4/3$. Take $\zeta,R$ from the dense gap calculation in
Section~\ref{subsec:ugc-tw-gap} for factor $r'$, and use
Theorem~\ref{thm:bounded-degree-kr-source} with $\eta=\zeta/2$,
as in Theorem~\ref{thm:ugc-tripartite-cover}. This is admissible
since $0<\eta<1/24<1/10$. Apply it in
$K_R(P_3(G))$, where $n=|V(G)|$.
Replace only the three global cliques by bounded-degree expanders
and retain every product edge in each source copy. These source
edges already have bounded degree, so
$d_1=1$ and $d_0=2d_\eta=D_{3,\zeta}$, with no biclique factors.
Lemma~\ref{lem:multipartite-product} gives YES cover bound
$(3/2+3\zeta/2)n\leq(3/2+3\zeta)n$ and NO cover bound
$(2-\zeta)n\geq(2-3\zeta)n$ for $P_3(G)$, as required by
Section~\ref{subsec:ugc-tw-gap}.
The global cliques give YES scale at least $Rn-1$, compared with
order $3Rn$. A small fixed $\tau$ therefore absorbs the loss into
the factor slack; see Appendix~\ref{app:transfer-tw-ugc}.

\emph{SSE hardness.}
Use the common line-graph factors.
Lemma~\ref{lem:treewidth-balanced-separator} gives
$b(H)\leq\operatorname{tw}(H)+1$ for the dense line graph.
Thus the sparse NO width is at least $\beta N-1-\tau N$.
The YES line-graph construction has treedepth at most $4\eta N$.
An elimination forest of height $h$ gives a tree decomposition of
width at most $h-1$: connect its roots and use its ancestor paths as
bags. A fixed vertex occurs on its descendant subtree, and every
edge is covered at its deeper endpoint. Hence the dense YES
treewidth is at most $4\eta N$. By
Lemma~\ref{lem:minor-monotonicity}, deleting edges to obtain the sparse
graph cannot increase this bound.
Choose $\eta$ for positive linear gap slack, then $\tau$ small and
$N$ large enough to absorb the sparsification loss and the constant
$1$. Appendix~\ref{app:transfer-tw-sse} gives the choices.

\end{proof}

The SSE treewidth hardness at every fixed factor on unrestricted
graphs is due to Wu, Austrin, Pitassi, and Liu~\cite[Theorem~1.1]{APW}. The bounded-degree
statement here follows by applying our treewidth sparsification
theorem to the line-graph construction.

For the total-depth transfer below, we use the following bound for
lifting an approximate solution through a reduction.

\begin{lemma}[Lifting approximate solutions]
\label{lem:solution-lifting}
Suppose each solution of a minimization instance $I'$ with cost $z$
can be converted in polynomial time into a solution of $I$ with cost
at most $z/a+\eta L(I)$, where $a>0$, $\eta\ge0$, and the bound $L(I)$ satisfies
\begin{equation}\label{eq:solution-lifting-lower}
 0\le L(I)\le\OPT(I).
\end{equation}
If
\begin{equation}\label{eq:solution-lifting-upper}
 \OPT(I')\le a\OPT(I)+\eta aL(I),
\end{equation}
then, for every $r\ge1$, an $r$-approximation for
$I'$ gives an $(r+(r+1)\eta)$-approximation for $I$. If instead
\begin{equation}\label{eq:solution-lifting-exact-upper}
 \OPT(I')\le a\OPT(I),
\end{equation}
the factor is $r+\eta$.
\end{lemma}

\begin{proof}
An $r$-approximate solution has $z\le r\OPT(I')$, so its conversion
has cost at most
\[
\frac za+\eta L(I)
 \le\frac{r\OPT(I')}{a}+\eta L(I) \overset{\eqref{eq:solution-lifting-upper}}{\le}
 r\OPT(I)+(r+1)\eta L(I)
 \overset{\eqref{eq:solution-lifting-lower}}{\le}
 \bigl(r+(r+1)\eta\bigr)\OPT(I).
\]
Under~\eqref{eq:solution-lifting-exact-upper}, the conversion cost satisfies
\[
 \frac za+\eta L(I)
 \overset{\eqref{eq:solution-lifting-exact-upper}}{\le}r\OPT(I)+\eta L(I)
 \overset{\eqref{eq:solution-lifting-lower}}{\le}(r+\eta)\OPT(I).
\]
\end{proof}

\Needspace{8\baselineskip}
\begin{corollary}[Bounded-degree total-depth hardness]
\label{cor:total-depth-bounded-degree}
\label{thm:np-total-depth-bounded-degree}
\label{thm:total-depth-ugc-bounded-degree}
\label{cor:total-depth-sse-bounded-degree}
Approximating total depth within every fixed factor $1\le r<63/62$
is NP-hard on simple unweighted graphs of bounded maximum degree.
On such graphs,
total depth is also Unique Games-hard to approximate within every
fixed factor $1\le r<8/7$, and SSE-hard to approximate within every
fixed factor $r\ge1$. The degree bound depends only on the excluded
factor. All reductions are deterministic.
\end{corollary}

\begin{proof}
Apply Theorem~\ref{thm:total-depth-sparsification}.
It loses at most $\tau N^2$ and completes sparse forests at the
same additive cost.

\emph{NP-hardness.}
Fix $1\le r<r'<63/62$. Use the dense family
$H=H^\Sigma_M(G_F^\Sigma)$ from
Theorem~\ref{thm:np-total-depth-dense-gap} for factor $r'$, where
$F$ is a parity system with $m$ equations and occurrence bound
$b_\epsilon$ from Theorem~\ref{thm:xor-source}.
The three part cliques and the three bicliques $X\times T_X$
for $X\in\{A,B,C\}$
give $d_1=2$. Retain every edge of $G_F^\Sigma$; its maximum degree
is at most $14b_\epsilon$ by~\eqref{eq:np-total-depth-degree}, so
$d_0=14b_\epsilon$. The part sizes are $27m,36m,36m$ by
\eqref{eq:np-total-depth-part-sizes}, while each $T_X$ has size
$M=\Theta(m)$. Thus all biclique side-size ratios are bounded by
a constant depending only on $r'$.

The dense family has $N=99m+3M=\Theta(m)$ and total depth at
least $Mm=\Omega(N^2)$. Choose $\tau$ small relative to this
quadratic scale and the gap slack. This preserves the gap and also
ensures $\totd(H')\ge Mm/2$ for the sparse graph $H'$, as proved
in Appendix~\ref{app:transfer-total-np}. The latter bound will be
used for regularization.

\emph{UGC hardness.}
Fix a target factor $r<8/7$, and choose $r'$ with
$r<r'<8/7$. Use the parameters $q,\eta$ from the dense gap
calculation in Section~\ref{subsec:ugc-total-depth-gap} for factor $r'$, starting from
Theorem~\ref{thm:bounded-degree-kr-source}. Replace the $q$ layer
cliques by bounded-degree expanders and retain every cross-layer
source edge. Thus $d_1=1$ and $d_0=(q-1)d_\eta$.
Write $n=|V(G)|$ and $N=qn$. The layer cliques force total depth at least
$q n(n+1)/2\ge N^2/(2q)$. Hence the loss $\tau N^2$ is a
small fraction of the YES scale when $\tau$ is small. See
Appendix~\ref{app:transfer-total-ugc} for the factor calculation.

\emph{SSE hardness.}
The YES line-graph construction has height at most $4\eta N$ and
has $N$ vertices, so its total depth is at most
$N\cdot4\eta N=4\eta N^2$. Lemma~\ref{lem:balanced-by-td} gives NO total depth
at least $\beta N^2/2$. The ratio of these quadratic bounds is unbounded as $\eta$ tends
to zero. Choose the dense gap first and then $\tau$ small enough
to absorb the quadratic loss; Appendix~\ref{app:transfer-total-sse}
gives the explicit slack.

\end{proof}

\subsection{Regular graphs and subdivisions}

Both constructions below recover a forest by contracting connected
vertex sets. The following observation records the common argument.

\begin{lemma}[Projection through connected parts]
\label{lem:connected-part-projection}
Let $G,H$ be finite simple graphs, and let
$(X_v)_{v\in V(G)}$ partition $V(H)$ into nonempty connected sets.
Suppose every edge $uv\in E(G)$ has an edge of $H$ between $X_u$
and $X_v$. From any elimination forest $F$ of $H$, one can compute
in polynomial time an elimination forest $\overline F$ of $G$ such that
\[
 \operatorname{depth}_{\overline F}(v)\le\operatorname{depth}_F(x)
 \quad(v\in V(G),\ x\in X_v).
\]
Consequently $h(\overline F)\le h(F)$ and
$\totd(F)\ge\sum_{v\in V(G)}|X_v|\operatorname{depth}_{\overline F}(v)$.
\end{lemma}

\begin{proof}
For each $v$, choose the shallowest vertex $r_v$ of $X_v$.
Lemma~\ref{lem:connected-set-ancestor} makes it unique and an ancestor
of every vertex in $X_v$. If $xy$ witnesses an edge $uv\in E(G)$,
with $x\in X_u$ and $y\in X_v$, then $x,y$ are comparable in $F$.
Both $r_u,r_v$ are ancestors of their deeper endpoint, so they too
are comparable. Restrict $F$ to the representatives and rename
$r_v$ as $v$; this gives $\overline F$. Restriction and the ancestor
property give, for every $x\in X_v$,
\[
 \operatorname{depth}_{\overline F}(v)
 \le\operatorname{depth}_F(r_v)
 \overset{\text{Lem.~\ref{lem:connected-set-ancestor}}}{\le}
 \operatorname{depth}_F(x).
\]
Taking the maximum and summing over the partition prove the two
consequences. Computing depths, choosing representatives, and restricting
the forest are polynomial-time operations. No connectedness of $G$ or
$H$ is required; for empty graphs the empty partition and forest suffice.
\end{proof}

\begin{proposition}\label{prop:regularization}
For every finite simple graph $G$ of order $n$ and fixed odd $d\geq3$ with
$d>\Delta(G)$, one can construct in polynomial time a simple
$d$-regular graph $R$ containing $G$ as an induced subgraph, with
$(d+2)n\leq |V(R)|\leq(d+3)n$ and
\[
\begin{gathered}
 \td(G)\leq\td(R)\leq\td(G)+d+2,\qquad \operatorname{tw}(G)\leq\operatorname{tw}(R)
 \leq\max\{\operatorname{tw}(G),d+2\},\\
 (d+2)\totd(G)\leq\totd(R)
 \leq(d+3)\totd(G)+\tfrac12(d+2)(d+3)n.
\end{gathered}
\]
\end{proposition}

\begin{proof}
For each $v\in V(G)$, set $t=d-\deg_G(v)$, so $1\le t\le d$.
Take graphs $Q_v$ whose vertex sets are pairwise disjoint and
disjoint from $V(G)$. Choose $P_v\subseteq V(Q_v)$ with $|P_v|=t$,
and join $v$ to every vertex of $P_v$. Define $Q_v$ as follows.
\begin{enumerate}
 \item If $t$ is even, take $K_{d+1}$ and remove a perfect matching
   on $P_v$.
 \item If $t\ge3$ is odd, take $K_{d+2}$ and remove a cycle with
   vertex set $P_v$ and a perfect matching on $V(Q_v)\setminus P_v$.
   The latter set has even size $d+2-t$.
 \item If $t=1$, write $P_v=\{p\}$, take $K_{d+2}$, and choose
   distinct $a,b\in V(Q_v)\setminus\{p\}$. Remove $pa,pb$ and a
   perfect matching on $V(Q_v)\setminus\{p,a,b\}$.
\end{enumerate}
Define the graph $R$ by
\[
V(R)=V(G)\mathbin{\dot\cup}
       \mathop{\dot\bigcup}_{v\in V(G)}V(Q_v),\qquad E(R)=E(G)\cup\bigcup_{v\in V(G)}E(Q_v)
       \cup\bigcup_{v\in V(G)}\{\{v,p\}:p\in P_v\}.
\]
The three cases give
\[
 \deg_{Q_v}(p)=d-1\quad(p\in P_v),\qquad
 \deg_{Q_v}(x)=d\quad(x\in V(Q_v)\setminus P_v).
\]
Hence $\deg_R(v)=\deg_G(v)+|P_v|=d$ for every $v\in V(G)$,
and $\deg_R(x)=d$ for every $x\in V(Q_v)$.
Each $Q_v$ has minimum degree $d-1$, and its order satisfies
\begin{equation}\label{eq:regularization-attachment-size}
 d+1\le |V(Q_v)|\le d+2\qquad(v\in V(G)).
\end{equation}
Every connected component therefore has at least $d$ vertices,
so $Q_v$ is connected because
$2d>d+2$ for $d\geq3$. The resulting graph $R$ is simple.

Given an elimination forest $F$ of $G$, appending $V(Q_v)$ as a chain
below $v$ increases the height by at
most $d+2$. If the chain has $|V(Q_v)|$ vertices, their total depth is
at most
\[
 |V(Q_v)|\operatorname{depth}_F(v)
 +\frac{|V(Q_v)|(|V(Q_v)|+1)}2
 \overset{\eqref{eq:regularization-attachment-size}}{\leq}
 (d+2)\operatorname{depth}_F(v)+\frac{(d+2)(d+3)}2.
\]
Summing over $v$ and including the original vertices gives the stated
total-depth upper bound. Attach a bag $\{v\}\cup V(Q_v)$ to an
existing bag containing $v$ in a tree decomposition of $G$.
Doing this for every $v$ proves the treewidth bound.

For the total-depth lower bound, partition $V(R)$ into connected sets
$X_v=\{v\}\cup V(Q_v)$. Each edge $uv$ of $G$ itself witnesses
an edge between $X_u$ and $X_v$. Thus any elimination forest $F$
of $R$ projects by Lemma~\ref{lem:connected-part-projection} to a
forest $\overline F$ of $G$ with
\begin{equation}\label{eq:regularization-depth-projection}
 \operatorname{depth}_F(x)
 \overset{\text{Lem.~\ref{lem:connected-part-projection}}}{\ge}
 \operatorname{depth}_{\overline F}(v)
 \qquad(v\in V(G),\ x\in X_v).
\end{equation}
Consequently,
\[
 \totd(F)
 \overset{\eqref{eq:regularization-depth-projection}}{\geq}
 \sum_{v\in V(G)}|X_v|\operatorname{depth}_{\overline F}(v)
 \overset{\eqref{eq:regularization-attachment-size}}{\geq}
 (d+2)\totd(\overline F)\geq(d+2)\totd(G).
\]
This argument also gives a polynomial-time projection of forests.
\end{proof}

\begin{corollary}\label{cor:regular-hardness}
The hardness ranges in Corollaries~\ref{cor:td-bounded-degree},
\ref{cor:tw-bounded-degree}, and~\ref{cor:total-depth-bounded-degree}
also hold on $d$-regular graphs for a suitable fixed odd degree $d$
depending only on the excluded factor.
\end{corollary}

\begin{proof}
Choose a source gap larger than the target gap and apply
Proposition~\ref{prop:regularization}. Its fixed additive costs for
treedepth and treewidth vanish relative to their linear YES scale.
For total depth, let $G$ be the source graph, put $n=|V(G)|$, and
divide costs by $d+2$: the error is
$O(n^2/d+dn)$. Choose $d$ large, then $n$ large, relative to the
quadratic scale. In the NP case, with $M,m$ from the NP construction,
the bound $L=Mm/2\le\totd(G)$
from Appendix~\ref{app:transfer-total-np} and the proposition's
forest projection give the same approximation transfer directly.
Appendix~\ref{app:transfer-regular} supplies all parameter choices
and the projection calculation.
\end{proof}

\begin{corollary}\label{cor:subdivision-hardness}
For every fixed $g\geq3$, the treedepth and treewidth hardness ranges above
hold on bounded-degree graphs that are bipartite, $2$-degenerate, and
have girth at least $g$. Here $2$-degenerate means that every
nonempty subgraph has a vertex of degree at most two, and girth is the
length of a shortest cycle (infinite for an acyclic graph). On the same class, total depth is
SSE-hard to approximate within every fixed factor.
\end{corollary}

\begin{proof}
Replace every edge of $G$ by a path of fixed even length
$\ell\geq\max\{2,g/3\}$, obtaining $G_\ell$. All original vertices
lie in one side of a bipartition, since every replacement path has
even length. Each internal vertex has degree two, and deleting all
internal vertices leaves the original vertices isolated, so
$G_\ell$ is $2$-degenerate. Original vertices retain their degrees,
so $\Delta(G_\ell)\leq\max\{\Delta(G),2\}$.
Every cycle uses at least three
replacement paths and therefore has length at least $3\ell\geq g$.
By Lemma~\ref{lem:minor-monotonicity}, contracting each replacement
path gives
$\operatorname{tw}(G)
\overset{\text{Lem.~\ref{lem:minor-monotonicity}}}{\leq}
\operatorname{tw}(G_\ell)$ and
$\td(G)\overset{\text{Lem.~\ref{lem:minor-monotonicity}}}{\leq}
\td(G_\ell)$. Conversely, for each path
$u=x_0,x_1,\ldots,x_\ell=v$, attach the bags
$\{u,v,x_1\},\{v,x_1,x_2\},\ldots,
\{v,x_{\ell-2},x_{\ell-1}\}$ along a path to an old bag containing
$u,v$. These bags cover the replacement path and preserve the
connectedness condition for each vertex, so
$\operatorname{tw}(G_\ell)\leq\max\{\operatorname{tw}(G),2\}$.
If $\operatorname{tw}(G)\ge2$, the two bounds give equality.
If $\operatorname{tw}(G)\le1$, then $G$ is a forest. Subdivision
preserves acyclicity and whether an edge exists, so it preserves
treewidth zero or one as well. Thus treewidth is preserved in every case.
A path on $s$ vertices has an
elimination tree of height at most $\lceil\log_2(s+1)\rceil$: choose
a middle vertex as root and recurse on the two remaining subpaths.
Thus the $\ell-1$ internal vertices on each replaced edge have such
a tree of height $h=\lceil\log_2\ell\rceil$. Attach this tree below
the deeper of the original edge's endpoints in an elimination forest
of $G$. The original endpoints are comparable, so the edge to the
shallower endpoint is represented as well. This gives
\begin{equation}\label{eq:subdivision-height-upper}
 \td(G_\ell)\leq\td(G)+h,\qquad h=\lceil\log_2\ell\rceil.
\end{equation}
Lemma~\ref{lem:gap-preservation} absorbs the fixed height increase
in~\eqref{eq:subdivision-height-upper}, since the YES bounds grow
linearly with the source order. For
total depth, start with a
$d$-regular instance $G$ on $n$ vertices from
Corollary~\ref{cor:regular-hardness}.
The identity $dn=2|E(G)|$ and oddness of $d$ imply that $n$ is even.
Add a perfect matching on the vertex set to make every degree even,
using auxiliary matching edges that may be parallel to original edges.
Each connected component of the augmented multigraph has even
degrees and therefore has an Euler tour, a closed walk traversing
every edge once. Orient its edges along this tour, then delete the
matching edges. Each vertex has outdegree
$(d+1)/2$ before deletion, so its outdegree in the original graph is
at least $(d-1)/2$.

Assign the internal vertices of each replacement path to the tail
of its oriented edge. For each vertex $v$, let $X_v$ consist of $v$
and the internal vertices assigned to it. The sets $X_v$ partition
$V(G_\ell)$. Each is connected, and
\begin{equation}\label{eq:subdivision-block-size}
 s=1+(\ell-1)(d-1)/2,\qquad |X_v|\ge s\quad(v\in V(G)).
\end{equation}
Each oriented edge $uv$ has a replacement path whose last edge
joins $X_u$ to $X_v$. Lemma~\ref{lem:connected-part-projection}
therefore projects any elimination forest $F$ of $G_\ell$ to an
elimination forest $\overline F$ of $G$, and gives
\begin{equation}\label{eq:subdivision-depth-projection}
 \totd(F)
 \overset{\text{Lem.~\ref{lem:connected-part-projection}}}{\ge}
 \sum_{v\in V(G)}|X_v|\operatorname{depth}_{\overline F}(v)
 \overset{\eqref{eq:subdivision-block-size}}{\ge}
 s\totd(\overline F)\ge s\totd(G).
\end{equation}

For the upper bound, use the same midpoint-recursive trees below the
deeper endpoints in an optimal elimination forest $F$ of $G$, and call
the resulting forest of $G_\ell$ $\widehat F$. Each of the
$\ell-1$ new vertices $x$ on a replacement path $uv$ satisfies
\begin{equation}\label{eq:subdivision-new-depth}
 \operatorname{depth}_{\widehat F}(x)\le
 \operatorname{depth}_F(\text{deeper endpoint of }uv)+h.
\end{equation}
Each vertex is the deeper endpoint of at
most $d$ edges, so summing the original and new depths gives
\[
\begin{aligned}
 s\totd(G)
 &\overset{\eqref{eq:subdivision-depth-projection}}{\leq}
 \totd(G_\ell)\\
 &\overset{\eqref{eq:subdivision-new-depth}}{\leq}
 \totd(G)+(\ell-1)\sum_{uv\in E(G)}
       \operatorname{depth}_{F}(\text{deeper endpoint of }uv)
       +h(\ell-1)|E(G)|\\
 &\leq
 [1+d(\ell-1)]\totd(G)+h(\ell-1)|E(G)|.
\end{aligned}
\]
Since $d\geq3$, using the definition of $s$, $|E(G)|=dn/2$,
and $C=h(\ell-1)d/2$ gives
\[
 3s-[1+d(\ell-1)]=2+\frac{(d-3)(\ell-1)}2\geq0,\qquad
 s\totd(G)\leq\totd(G_\ell)\leq3s\totd(G)+Cn.
\]
Choose a regular SSE source with a gap larger than three times
the target factor. This fixes $d$; next fix $\ell$, and then enlarge
the source so that $Cn$ is negligible compared with its quadratic
YES scale. The resulting gap survives by the calculation in
Appendix~\ref{app:transfer-subdivision}.

\end{proof}

Equation~\eqref{eq:tp-adjusted-cost} also transfers the total-depth
hardness conclusions of both corollaries to trivially perfect
completion.
